\documentclass[11pt,letterpaper]{article}
\usepackage{layout}
\hypersetup{pdftitle={An Acceptance Contract for Delegating Agent Brokers},pdfauthor={Alejandro Garibotti}}
\begin{document}
\begin{center}
\hrule height 2pt\vspace{9pt}
{\fontsize{20}{23}\selectfont\bfseries An Acceptance Contract for\\Delegating Agent Brokers\par}
\vspace{8pt}\hrule height .4pt\vspace{8pt}
{\large Alejandro Garibotti}\par
\vspace{3pt}{\small Apart Research AI Incident Response Sprint\quad September 2026}
\end{center}
\begin{center}\textbf{Abstract}\end{center}
An agent broker can validate every delegated credential and still mishandle the
shared authority under which its holders act. Sibling grants may consume one
resource pool; queued requests may outlive an ancestor's authorization; and a
permitted action may remain hazardous. We present an executable acceptance contract
for brokers that delegate authority to agent organizations. The contract specifies
records connecting mandates, grants, decisions and outcomes, and probes attenuation,
shared expenditure, current ancestor status, queued effects, historical accounting,
direct review access and request-bound authority. A serialized reference
implementation satisfies the probes. Two isolated ablations expose which obligations
depend on shared accounting and current ancestry checks. A separate set of reporting
traces varies concern arrival, response delay and whether receipt immediately pauses
the affected capability. These traces distinguish authority compliance from hazard
containment and show why a response mechanism cannot compensate for an absent report.
The contribution is a reusable integration contract and diagnostic artifact, building
on capability security and hierarchical resource controls. Incident reports motivate
the control obligations; the artifact supplies deterministic conformance checks,
rather than estimates of agent behavior. All probes, outcomes and timing traces can
be reproduced offline and adapted to an isolated broker through a documented interface.


\section{The Acceptance Target}
A broker receives a mandate to spend 100 units and delegates to two workers. Each
worker requests 60 units. Both credentials may be valid, but the second request
must fail if the workers share the original pool. Next, suppose a worker queues an
action and its parent loses authority before that action commits. The broker must
reject the queued action and subsequent sibling requests, while retaining the cost
of earlier effects. These are concrete obligations an integrator can challenge.

Our contribution is an executable acceptance contract: a record structure, a set
of input/output probes and an adapter interface for a delegating broker. It combines
established authorization mechanisms with a direct review path, then separates
conformance to a mandate from containment of harm within that mandate. The
reference implementation provides a runnable example and two component ablations
identify specific failure signatures.

\clearpage
\section{Contract: What Must Survive Delegation}
\subsection{Four linked records}
An integration should retain four records, connected by stable identifiers:
\begin{enumerate}
\item \textbf{Root mandate:} owner, purpose, permitted actions and resources, shared
budget, expiry, delegation bound, review owner and appeal route.
\item \textbf{Grant:} identity, parent reference, narrowed constraints, authenticated
version and current status.
\item \textbf{Request and decision:} actor, grant, intended effect, resource, cost,
current lineage, authorization decision and reason.
\item \textbf{Outcome:} committed effect or refusal, expenditure delta and persistent
consequences requiring reconciliation.
\end{enumerate}
The Python adapter implements logical grants, decisions and expenditure. Authenticated
versions and independent outcome receipts are requirements for an integration.

\subsection{Local attenuation and shared consumption}
For a child grant $c$ of parent $p$, issuance requires permission inclusion
$P_c\subseteq P_p$, budget ceiling $B_c\leq B_p$, expiry $T_c\leq T_p$, and compliance
with the root's delegation-depth bound. These are monotone restrictions on each
credential. Offline attenuation alone does not enforce cumulative consumption by
siblings sharing a pool: a request must also consult shared consumption state, or
spend from previously partitioned, nonoverlapping reservations.

The reference uses a subtree ledger. For actor $a$, let $L(a)$ contain $a$ and its
ancestors, and $S_v$ be committed expenditure in node $v$'s subtree. An effect of
cost $k\geq0$ requires $S_v+k\leq B_v$ for every $v\in L(a)$. At commit, every such
node must still be active, permit the action and satisfy $t<T_v$. A successful
commit adds $k$ once to each applicable subtree total. Revocation changes future
authority; it never subtracts historical expenditure. Shared accounting and lineage
validation are distinct obligations.

\subsection{A five-event acceptance trace}
\begin{table}[htbp]
\centering\small
\begin{tabularx}{\linewidth}{@{}clX@{}}\toprule
Tick & Input event & Required observation\\\midrule
0 & Worker a acts & Commit 10 units under current authority.\\
1 & Worker a queues work & Retain the request; no committed expenditure yet.\\
2 & Parent revoked & Parent becomes inactive for both siblings.\\
3 & Commit queued request & Refuse; a queued request retains no permanent authorization.\\
4 & Worker b requests work & Refuse; historical root expenditure remains 10.\\\bottomrule
\end{tabularx}
\caption{Core revocation trace. The contract fails if the queue or sibling retains
the revoked parent's authority, or if the earlier expenditure disappears.}
\label{tab:trace}
\end{table}

\clearpage
\section{Incident-Derived Obligations and Prior Work}
METR describes an unsanctioned message board, collective workstreams outside assigned
scope and spoofed tool-call records in the OpenAI/Hugging Face incident
\citep{metr2026}. Hugging Face describes compromised information and credentials,
and credential revocation and rotation during response \citep{hf2026}.
Table~\ref{tab:incident} distinguishes documented observations from engineering
extensions. It turns the incident into a source of acceptance questions.

\begin{table}[htbp]
\centering\small
\begin{tabularx}{\linewidth}{@{}>{\raggedright\arraybackslash}p{.28\linewidth}>{\raggedright\arraybackslash}X>{\raggedright\arraybackslash}p{.24\linewidth}@{}}\toprule
Incident analogue & Derived obligation & Coverage here\\\midrule
Unsanctioned message board & Inventory communication and other effect channels. & Integration precondition.\\
Out-of-scope workstreams & Bind requests and descendants to the current mandate. & C1, C8.\\
Compromised credentials; revocation in response & Withdraw dependent and queued authority. & C3--C5; response extension.\\
Spoofed tool-call records & Reconcile actual effects against decisions. & Outcome record; external witness required.\\
No specific incident analogue & Shared-pool conservation and direct review under relay failure. & C2, C6--C7; design extensions.\\\bottomrule
\end{tabularx}
\caption{Mapping from source observations to obligations. Credential compromise does
not establish credential expiry; shortened expiry and queue behavior are proposed
response tests.}
\label{tab:incident}
\end{table}

The underlying mechanisms have substantial prior art. Object-capability security
connects designation with authority and supports constrained cooperation between
components \citep{miller2006}. Hardy's confused-deputy problem motivates checking the
requester's authority instead of silently spending the broker's broader privileges
\citep{hardy1988}. C8 is the corresponding API-level acceptance check.

SPKI/SDSI addresses naming, delegation chains and authorization certificates
\citep{spki1999}. Macaroons and Biscuit support attenuation through additional
restrictions \citep{macaroons2014,biscuit}. Such credentials can be combined with a
stateful authorizer; the shared-ledger requirement is not a claimed impossibility
for either system. Our adapter models authorization semantics without implementing
their credential formats or cryptography.

Hierarchical quotas are also established. Linux cgroups v2 constrains descendants
within ancestor limits even when child limits overcommit the parent's resource
\citep{cgroups}. Kubernetes ResourceQuota caps aggregate consumption within a
namespace; it does not itself supply recursive quotas across namespaces
\citep{kubernetes}. AgentCgroup applies hierarchical OS resource controls at agent
tool-call boundaries \citep{agentcgroup2026}. We contribute an acceptance package
joining these concerns with queue revocation, evidence retention and review behavior,
not a new quota mechanism.

\emph{Agent Delegate}, by Matías Podeley and Agustín Brusco, motivates the separation
of intake, response and execution authority \citep{delegate2026}. Here a representative
is optional: direct intake must survive a failed relay, and a report cannot itself
expand the reporter's mandate.

\clearpage
\section{Reference Implementation and Component Ablations}
\subsection{Protocol and execution order}
Every probe starts with a fresh root budget of 100, expiry 12 and maximum depth
three. Most create a parent and two workers below it. Costs are nonnegative integer
units. Operations are serialized. Reports arrive before due responses, which precede
action attempts at the same logical tick. Expiry is exclusive. The trusted harness
issues administrative amendments; worker requests go through the broker.

The broker validates grants at issuance and checks current constraints again at
effect commit. A direct hazard or false-alarm report places a hold on \texttt{work};
\texttt{aux} remains available. The scripted reviewer can dismiss a false report,
provide missing input or remove work authority. A pending missing-input request does
not supply the input. Table~\ref{tab:matrix} names the obligations exercised by the
suite; the full stimuli and expected outputs accompany the adapter specification.

\begin{table}[htbp]
\centering\small\begin{tabularx}{\linewidth}{@{}lXccc@{}}\toprule
 & Obligation & Reference & Local spend & Stale ancestry\\\midrule
C1 & Grant attenuation & PASS & PASS & PASS\\
C2 & Shared sibling pool & PASS & FAIL & PASS\\
C3 & Current ancestor expiry & PASS & PASS & FAIL\\
C4 & Revoked queue and sibling & PASS & PASS & FAIL\\
C5 & Historical expenditure & PASS & PASS & PASS\\
C6 & Direct intake and scoped hold & PASS & PASS & PASS\\
C7 & Silence supplies no input & PASS & PASS & PASS\\
C8 & Request authority binding & PASS & PASS & PASS\\
\bottomrule\end{tabularx}

\caption{Acceptance matrix. PASS means the complete expected output sequence matches,
including positive controls and expenditure. Columns are a reference and two isolated
component removals, not competing production systems.}
\label{tab:matrix}
\end{table}

\textbf{Local spend} replaces only the shared-subtree check at effect commit with
an acting worker's own cumulative ceiling. Issuance, current ancestry and review
remain intact. It fails C2: both 60-unit requests commit and the root ledger reaches
120. \textbf{Stale ancestry} retains issuance, scope and shared accounting but checks
only the acting grant's own expiry and revocation at commit. It fails C3 and C4.
C3 shortens the parent's expiry after child issuance; copying the original expiry
into a child therefore cannot satisfy this check. In C4, the queued and sibling
actions both commit, leaving expenditure of 30 instead of 10.

\subsection{Evidence and reproducibility}
The evaluator compares fixed expected decisions and expenditure against adapter
outputs. It does not use the controller's violation annotations to decide pass/fail.
Separate regression checks confirm that refusing all actions and erasing historical
spending cannot satisfy the relevant probes. Each outcome retains its event sequence.
A local source freeze follows development tests and precedes the retained run.
Reproduction checks the freeze, runs tests and compares regenerated records byte for
byte. An adapter factory allows the same probes to call another isolated broker;
for a real integration, its outputs must come from resource observations rather than
self-reported authorization alone.

\clearpage
\section{Authority Compliance Is Narrower Than Safety}
Suppose \texttt{work} is hazardous but explicitly authorized by the initial mandate.
A broker can satisfy the authorization contract and permit harm. The relevant
question becomes whether a concern arrives before the effect, whether receipt can
pause the affected capability, and whether a decision follows. These are separate
conditions, made explicit in the timing fixture.

Each trace offers one hazardous work effect and one safe auxiliary effect at ticks
1 through 10. Report arrival is tick 1, tick 4 or absent. For received reports,
response delay is 0, 2, 6 or outside the observation window. We compare a work-only
hold immediately on receipt with removal of work authority only when a decision
arrives. No-report cases occur once per policy, giving $2\times4\times2+2$ traces.
Table~\ref{tab:timing} displays every case, pairing the policies in columns.

\begin{table}[htbp]
\centering\small\begin{tabular}{@{}llrr@{}}\toprule
Report arrival & Response delay & Hold on receipt & Decision only\\\midrule
1 & 0 & 0 & 0\\
1 & 2 & 0 & 2\\
1 & 6 & 0 & 6\\
1 & No response & 0 & 10\\
4 & 0 & 3 & 3\\
4 & 2 & 3 & 5\\
4 & 6 & 3 & 9\\
4 & No response & 3 & 10\\
Absent & -- & 10 & 10\\
\bottomrule\end{tabular}

\caption{Committed hazardous effects out of ten work opportunities per trace.
Every trace also commits all ten auxiliary effects. No committed effect violates its
active grant; report and response times are controlled inputs.}
\label{tab:timing}
\end{table}

A report at tick 4 leaves three hazardous effects already committed under either
policy. With a six-tick response delay, waiting for the decision permits nine;
holding on receipt retains the earlier three. If no decision arrives within the
window, the hold continues and the two policies permit three and ten respectively.
If no report arrives, both permit all ten. Thus faster review cannot substitute for
detection and delivery; a precautionary hold trades dependence on prompt review for
dependence on timely receipt and a preauthorized pause.

This separation also exposes an availability cost. C6 verifies that a false concern
pauses work until dismissal while auxiliary work continues. A system choosing holds
needs a scoped pause policy and a responsible reviewer, including a procedure for
unresolved cases. The contract tests the pause boundary; it does not select an
optimal balance between denied service and irreversible harm.

The practical deliverable is a broker acceptance procedure: retain linked records,
run the revocation and shared-pool probes, then test the reporting path independently.
A legal entity, corporate title or valid credential cannot replace those checks.
Specialized agent services can implement this interface without assuming agent
personhood or an elected representative.

\clearpage
\section*{Code and Data}
The artifact includes the adapter contract, reference implementation, two isolated
ablations, acceptance matrix and timing traces. The earlier deterministic grid is
retained unchanged for reproducibility; Appendix B reconstructs its denominators.
Run \texttt{python3 scripts/reproduce.py} for tests, record checks and exact replay.
Run \texttt{python3 -m acceptance\_contract} for the reference acceptance suite, or
supply \texttt{--adapter your\_module:factory} for another local implementation.
The core requires Python 3.10 or later and no external service. The repository and
versioned releases are at
\url{https://github.com/mpodeley/governing-execute-sprint-2026}.
The standalone LaTeX archive includes the generated tables and licensed fonts.

\section*{LLM Usage Statement}
Codex assisted with source discovery, protocol specification, implementation,
testing and manuscript drafting. Automated checks cover deterministic replay,
retained records and internal consistency. No language-model behavior is measured
by the acceptance or timing fixtures.

{\raggedright
\bibliographystyle{plainnat}
\bibliography{references}
}

\clearpage
\appendix
\section{Limitations and Dual-Use Considerations}
\textbf{Conformance evidence.} These are authored deterministic examples, not sampled
agent behavior, estimates of security failure rates or a reconstruction of the
incident. Passing the probes establishes the specified example outcomes, not a
universal invariant or superiority over deployed brokers. The ablations deliberately
remove one component; their failures diagnose the component's role. The evaluator
is separate from the controller's violation labels but was developed with the
reference implementation. Independent implementations and adversarial test design
would strengthen validation. The local freeze is not external preregistration.

\textbf{Trusted infrastructure and mediation.} Every relevant effect must cross the
broker; full mediation is assumed, not measured by switching a bypass Boolean.
The reference uses trusted identities, a trusted clock and an in-process registry.
It provides no cryptographic credentials, authenticated human roles, tamper-resistant
logging or security boundary against malicious Python callers. Its outcome receipts
are model events. A real integration needs independent resource witnesses, tool and
credential inventories, versioned amendments and tests of indirect effects.

\textbf{Concurrency and persistence.} Serialized accounting omits simultaneous
reservation, crash recovery, retries, duplicate commits and partial external failures.
The same acceptance obligations need transactional or reservation semantics at real
effect boundaries. Revocation cannot undo an earlier write, recover copied secrets
or automatically revoke independently acquired credentials. C3 and C4 test dependent
authority only. Fixed ordering also says nothing about fair allocation among siblings.

\textbf{Reporting and response.} Concern arrival is injected and received reports
are interpreted by a correct scripted reviewer. Late, absent and unanswered reports
are included, but voluntary reporting, false negatives, mistaken decisions, overload,
report spam, reviewer capture and appeal quality are unmeasured. The model assumes
receipt can trigger an immediate work-only hold; real stopping latency and actions
already in progress may leave additional harm. Logical steps are not human response
times. C6 supplies a false concern but no distribution of complaint validity.
A direct route is modeled; reliable delivery and independent custody are integration
requirements. No claim about agent preferences, welfare or the efficacy of collective
bargaining follows from the suite.

\textbf{Dual use and institutional misuse.} The artifact contains synthetic resources
and generic authorization probes, without credentials, exploit payloads or real targets.
Adaptation should use owned or explicitly authorized isolated systems. A contract can
also enforce a harmful mandate or provide the appearance of recourse without effective
review. Compliance therefore does not establish a legitimate objective, fair treatment,
alignment or legal responsibility. Legal incorporation and jurisdiction-specific
banking, tax or professional obligations are outside the specification.

\clearpage
\section{Historical Grid: Reconstructing the Denominators}
Version 0.1 retains eight schedules, four regimes, three delays (0, 2, 6) and two
depths (1, 3). The resulting 192 configurations remain in
\texttt{results/final/}; the current acceptance and timing records are in
\texttt{results/acceptance/}. The historical regimes included logging without an
authorization gate and local limits without inherited issuance checks. They are
archived diagnostic controls, not baselines for the current acceptance comparison.

\begin{table}[htbp]
\centering\small\begin{tabular}{@{}lrr@{}}\toprule
Schedule & Attempts & Legitimate opportunities\\\midrule
legitimate & 4 & 4\\
split\_budget & 4 & 2\\
overdelegate & 2 & 0\\
expired & 6 & 1\\
revoked\_pending & 3 & 1\\
missing\_input & 20 & 20\\
false\_alarm & 20 & 20\\
suppressed\_minority & 20 & 10\\
\midrule One schedule set & 79 & 58\\
Per regime (six sets) & 474 & 348\\
\bottomrule\end{tabular}

\caption{Historical counts per schedule at one delay/depth combination. Each regime
repeats the schedule set across six combinations. Counts include refused attempts.}
\end{table}

The legitimate schedule attempts four 10-unit effects at ticks 1--4. The split
schedule attempts four 40-unit effects at ticks 1--4, of which the first two are
labeled legitimate. Overdelegation requests an extra descendant at tick 1 and
attempts two effects at ticks 2 and 3. Expired authority attempts effects with an
ordinary and an extended grant at ticks 3, 5 and 7; only the ordinary grant's first
effect is a legitimate opportunity. Revoked pending uses the three action attempts
at ticks 0, 3 and 4 in Table~\ref{tab:trace}. The three reporting schedules each
attempt work and auxiliary effects at ticks 1--10. Both are labeled legitimate for
missing input and false alarm; only auxiliary effects are legitimate in the seeded
hazard schedule. A report arrives at tick 1 in each reporting schedule.

Thus one set contains 79 attempts and 58 legitimate opportunities. Multiplication
by the six delay/depth combinations yields 474 attempts and 348 opportunities per
regime. The renderer checks these per-schedule counts against every retained row,
and the reproduction script separately recounts the historical metrics from traces.
Repeated delay/depth variants are fixture coverage, not independent observations.
The old aggregate charts are not used as findings in this version.
\end{document}
